\chapter{Manual de uso: Gemini en las tareas de la tutoría personalizada}
\label{anexo:manual_ia}

Este manual es un producto del segundo objetivo específico (Sección~\ref{sec:metodologia_por_objetivo}). Está dirigido a los estudiantes que participan en la secuencia de tareas y explica cómo usar Gemini dentro de ella. No es un manual general de inteligencia artificial: es una guía para este diseño.

\section*{¿Qué es Gemini?}

Gemini es un asistente de inteligencia artificial de Google con el que se puede conversar por escrito, en español, para hacer preguntas, pedir explicaciones o discutir una idea. No es una calculadora ni un buscador. Genera sus respuestas a partir de un modelo de lenguaje, y eso significa dos cosas: puede explicar muy bien un concepto y, en otras ocasiones, equivocarse con total seguridad. En matemáticas, los errores más comunes son de cálculo y de pasos que parecen lógicos pero no lo son.

\section*{¿Cómo se usa en la tutoría?}

\begin{itemize}
\item Se usa en el computador de la tutoría, con la cuenta institucional del tutor. Usted no necesita su cuenta personal.
\item Se usa con el modo \emph{aprendizaje guiado} activado. En ese modo, Gemini responde con preguntas, pistas y explicaciones por partes, en lugar de dar la solución completa. Si le hace una pregunta y le responde con otra pregunta, no es un error: así funciona este modo.
\item No escriba en la conversación su nombre, su documento ni otros datos personales.
\item Al terminar la sesión, la conversación se guarda sin datos que lo identifiquen y se borra de la cuenta.
\end{itemize}

\section*{¿Cuándo se usa?}

Cada tarea tiene tres momentos. Gemini solo se usa en el segundo:

\begin{description}
\item[Momento 1 --- sin IA:] usted trabaja solo, con esfero azul o negro. Aquí construye su propia respuesta, la que después va a comparar con lo que diga la IA.
\item[Momento 2 --- con Gemini:] puede consultar lo que quiera. Si decide cambiar algo en su hoja, hágalo con esfero rojo, sin borrar lo que había escrito.
\item[Momento 3 --- validación:] compara lo que hizo en los dos momentos y escribe su respuesta final con su justificación.
\end{description}

\section*{Cómo formular una buena pregunta}

La diferencia entre una consulta que le ayuda a pensar y una que solo le da la respuesta está en cómo pregunta.

\begin{table}[H]
\centering
\caption{Ejemplos de preguntas a Gemini: qué evitar y qué preguntar en su lugar}
\small
\begin{tabular}{|p{6.8cm}|p{6.8cm}|}
\hline
\textbf{Evitar (pide la respuesta)} & \textbf{Preguntar en su lugar (pide ayuda para pensar)} \\
\hline
``¿Cuánto da el límite de $\frac{x^2-4}{x-2}$ cuando $x$ tiende a 2?'' &
``Reemplacé $x$ por 2 y me dio $0/0$. ¿Qué significa eso y qué podría intentar?'' \\
\hline
``¿Este razonamiento está bien?'' &
``Pienso que este razonamiento está mal porque [su argumento]. ¿Mi argumento tiene sentido?'' \\
\hline
``Resuélvame el ejercicio.'' &
``Llegué hasta aquí: [sus pasos]. ¿En qué paso cree que me equivoqué? No me diga la respuesta.'' \\
\hline
\end{tabular}
\end{table}

\section*{Qué hacer con la respuesta}

\begin{enumerate}
\item \textbf{Compárela con lo que usted hizo en el momento 1.} Si no coincide, esa diferencia es importante: pregúntese cuál de los dos tiene razón y por qué.
\item \textbf{Verifíquela.} Puede reemplazar valores, hacer una tabla o una gráfica en GeoGebra. Que Gemini lo diga con seguridad no significa que sea correcto.
\item \textbf{Pida otra explicación si no entendió.} Con un ejemplo numérico, con una gráfica o con otras palabras.
\item \textbf{No copie lo que no puede explicar.} Si va a usar una idea de Gemini en su respuesta, debe poder decir por qué es correcta.
\end{enumerate}

\section*{Para el tutor: cómo introducir la herramienta}

En la sesión de orientación (Fase 1, Sección~\ref{sec:metodos_materiales}) se presenta este manual y se hace una demostración breve: una misma duda formulada de dos maneras ---pidiendo la respuesta y pidiendo ayuda para pensar--- y lo que Gemini responde en cada caso. Si el estudiante asistió al taller de Gemini del programa GEA, se retoma lo que aprendió allí. No se trata de una capacitación técnica, sino de dejar claro para qué se usa la herramienta en las tareas antes de la primera consulta real.

\section*{Privacidad y registro}

Las conversaciones con Gemini durante las tareas se registran como datos del estudio (Sección~\ref{sec:consideraciones_eticas}): la pregunta, la respuesta y lo que el estudiante decide hacer con ella. El estudiante lo sabe desde el consentimiento informado. Ninguna conversación se publica ni se comparte fuera de los fines del estudio.
